\documentclass[10pt,a4paper]{amsart}
\usepackage[margin=19mm]{geometry}
\usepackage{amsmath,amssymb,mathtools}
\usepackage[scr=rsfs,cal=euler]{mathalfa}
\usepackage{xfrac}
\usepackage[unicode,hidelinks,bookmarks=false]{hyperref}
\newcommand{\ie}{i.e.\ }
\newcommand{\cf}{cf.\ }
\newcommand{\resp}{respectively\ }
\newcommand{\Subsection}{\subsection}
\providecommand{\nobreakdash}{\nobreak\hbox{-}}
\setlength{\emergencystretch}{2em}
\multlinegap0pt
\usepackage{amssymb}
\usepackage{xcolor}
\newtheorem{theorem}{Theorem}
\newtheorem{remark}{Remark}
\newtheorem{lemma}{Lemma}
\usepackage{cleveref}
\crefname{theorem}{Theorem}{Theorems}
\crefname{remark}{Remark}{Remarks}
\crefname{lemma}{Lemma}{Lemmas}
\newcommand{\TC}{\operatorname{TC}}
\newcommand{\Tr}{\operatorname{Tr}}
\newcommand{\bbC}{\mathbb{C}}
\newcommand{\bbN}{\mathbb{N}}
\title{Metrics on completely positive maps via noncommutative geometry}
\author{Are Austad, Erik B\'edos, Jonas Eidesen, Nadia S. Larsen, Tron Omland}
\date{Source: arXiv:2512.10842v3 (CC BY 4.0)}
\begin{document}
\maketitle

\noindent\emph{Source and licence.} Metrics on completely positive maps via noncommutative geometry. Version arXiv:2512.10842v3 (CC BY 4.0), distributed by arXiv under the Creative Commons Attribution 4.0 International licence, http://creativecommons.org/licenses/by/4.0/, as recorded in the arXiv licence metadata for this exact version and shown on the abstract page https://arxiv.org/abs/2512.10842v3. Full licence notice: \texttt{licenses/arxiv-2512.10842-CC-BY-4.0.txt}.\par\smallskip

\noindent\textit{Editorial scope.} Counted unit: the article's theorem 2 (source0.tex lines 1584--1591). The article's complete original proof of it is delivered with it (source0.tex lines 1592--1631, one \texttt{proof} environment of 40 lines). No mathematics is written, repaired or re-derived: the statement and the proof are the article's own lines, in the article's own notation, and no retained line is substituted or re-flowed.  Retained uncounted as the source-local context the proof needs: lines 1246--1253; lines 1277--1279; lines 1556--1562.  Nothing was omitted from any retained span, so the package carries no omission record.  No retained line contains a citation, so the wrapper carries no bibliography and no bibliography entry.  No retained line contains a figure environment, an image-inclusion command or drawing code, so no image is delivered and the package declares no asset dependency.  Editorial formatting only, in addition: an AMS \texttt{amsart} preamble that restates the article's own declarations and macros used by the retained lines, namely the environments \texttt{theorem}, \texttt{remark}, \texttt{lemma} on separate counters, together with the standard packages those lines need.  The retained environments print in this document as Theorem 1, Remark 1, Lemma 1 in order of appearance, whereas the article prints the same items with the numbering of its own full document.  The complete native source of the article as distributed in the arXiv e-print for this exact version is bundled unchanged as \texttt{sources/190840/source0.tex}.
\begin{theorem}\label{thm:stability 1}
    Let $A$ be a unital $C^*$-algebra, and $B$ be a $C^*$-algebra with a faithful trace $\tau$. For each $n \in \bbN$ suppose that we have a seminorm $L_n$ on $(M_n(\bbC) \otimes A) \otimes_{\rm max} (M_n(\bbC) \otimes B)^{\rm op}$ such that
    \begin{enumerate}
        \item $(1_{M_n(\bbC) \otimes M_n(\bbC)^{\rm op}} \otimes L_1) \circ \Sigma_{[23]} \leq L_n$, and
        \item $L_n(\Sigma_{[23]}(1_n \otimes 1_n^{\rm op} \otimes x)) \leq 1$ for every $x \in A \otimes_{\rm max} B^{\rm op}$ satisfying $L_1(x) \leq 1$.
    \end{enumerate}
    Then the sequence $\{\Delta_{\Tr_n \otimes \tau, L_n}^{\rm max}\}_{n \in \bbN}$ of extended metrics on $\TC_{\Tr_n \otimes \tau} (M_n (\bbC) \otimes A , M_n(\bbC) \otimes B)$ is stable.
\end{theorem}

\begin{remark}\label{rem:stability for spatial norm}
    If $B$ is a unital $C^*$-algebra with an amenable faithful trace $\tau$, then the above result and proof go through if we replace the maximal tensor product with the minimal tensor product everywhere.
\end{remark}

\begin{lemma}\label{thm:spectral triples Kaad}
    Let $(A,H_A,\partial_A)$ and $(B,H_B,\partial_B)$ be unital spectral triples. Then for any $x \in A \otimes_{\rm min} B$ we have that
    \begin{align*}
        (1_A \otimes L_{\partial_B})(x) & \leq L_{\partial_A \times \partial_B}(x), \text{ and} \\
        (L_{\partial_A} \otimes 1_B)(x) & \leq L_{\partial_A \times \partial_B}(x).
    \end{align*}
\end{lemma}

\begin{theorem}\label{thm:stability 2}
    Let $(A,H_A,\partial_A)$, and $(B^{\rm op},H_B,\partial_B)$ be unital spectral triples, and assume that $B$ admits an amenable faithful trace $\tau$. For each $n \in \bbN$, let $(M_n(\bbC),H_n,\partial_n)$ be a spectral triple. By identifying $M_n(\bbC)^{\rm op}$ with $M_n(\bbC)$ we may view $(M_n(\bbC),H_n,\partial_n)$ as a spectral triple for $M_n(\bbC)^{\rm op}$ as well. Let
    \begin{align*}
        L_1 & = L_{\partial_A \times \partial_B}, \text{ and } \\
        L_n & = L_{(\partial_n \times \partial_n) \times (\partial_A \times \partial_B)} \circ \Sigma_{[23]}, \text{ for } n \geq 2.
    \end{align*}
    Then the sequence $\{\Delta_{\Tr_n \otimes \tau, L_n}^{\rm min}\}_{n \in \bbN}$ of extended metrics on $\TC_{\Tr_n \otimes \tau}(M_n(\bbC) \otimes A, M_n(\bbC) \otimes B)$ is stable.
\end{theorem}

\begin{proof}
    For ease of notation, let $\partial = (\partial_n \times \partial_n) \times (\partial_A \times \partial_B)$. To apply \cref{thm:stability 1} and \cref{rem:stability for spatial norm} we need to check that the following two conditions are satisfied:
    \begin{enumerate}
        \item $(1_{M_n(\bbC) \otimes M_n(\bbC)^{\rm op}} \otimes L_1) \circ \Sigma_{[23]} \leq L_n$, and
        \item $L_n(\Sigma_{[23]}(1_n \otimes 1_n^{\rm op} \otimes x)) \leq 1$ for every $x \in A \otimes_{\rm min} B^{\rm op}$ satisfying $L_1(x) \leq 1$.
    \end{enumerate}
    By \cref{thm:spectral triples Kaad} we have that
    \begin{equation*}
        1_{M_n(\bbC) \otimes M_n(\bbC)^{\rm op}} \otimes L_{\partial_A \times \partial_B} \leq L_{\partial}.
    \end{equation*}
    Since pre-composing with $\Sigma_{[23]}$ does not change this inequality we get that
    \begin{equation*}
        (1_{M_n(\bbC) \otimes M_n(\bbC)^{\rm op}} \otimes L_1) \circ \Sigma_{[23]} \leq L_n.
    \end{equation*}
    It remains to show that $L_n(\Sigma_{[23]}(1_n \otimes 1_n^{\rm op} \otimes x)) \leq 1$ for every $x \in A \otimes_{\rm min} B^{\rm op}$ satisfying $L_1(x) \leq 1$. By the definition of $L_n$, this means that we want to show that
    \begin{equation*}
        L_{\partial}(1_n \otimes 1_n^{\rm op} \otimes x) \leq 1
    \end{equation*}
    for every $x \in A \otimes_{\rm min} B^{\rm op}$ satisfying $L_{\partial_A \times \partial_B}(x) \leq 1$. We will in fact prove the stronger statement that
    \begin{equation}\label{eqn:seminorm from external products}
        L_{\partial}(1_n \otimes 1_n^{\rm op} \otimes x) = L_{\partial_A \times \partial_B}(x)
    \end{equation}
    for all $x \in A \otimes_{\rm min} B^{\rm op}$. To prove this we only need to consider the parity of the spectral triple for $M_n(\bbC) \otimes M_n(\bbC)^{\rm op}$, and the parity of the spectral triple for $A \otimes_{\rm min} B^{\rm op}$. In fact, since we are considering the same spectral triple for $M_n(\bbC)$ and $M_n(\bbC)^{\rm op}$, the resulting external product of these two spectral triples will always be even. Hence, we only need to consider the parity of the spectral triple for $A \otimes_{\rm min} B^{\rm op}$. In both cases we see that
    \begin{equation*}
        \partial
        = (\partial_n \times \partial_n) \otimes 1
        + \gamma_{M_n(\bbC) \otimes M_n(\bbC)^{\rm op}} \otimes (\partial_A \times \partial_B).
    \end{equation*}
    Hence, for any $x \in A \otimes_{\rm min} B^{\rm op}$ we have that
    \begin{align*}
        L_\partial(1_n \otimes 1_n^{\rm op} \otimes x)
        & = \|[\partial, 1_n \otimes 1_n^{\rm op} \otimes x]\| \\
        & = \|[(\partial_n \times \partial_n) \otimes 1 + \gamma_{M_n(\bbC) \otimes M_n(\bbC)^{\rm op}} \otimes (\partial_A \times \partial_B), 1_n \otimes 1_n^{\rm op} \otimes x]\| \\
        & = \|[\gamma_{M_n(\bbC) \otimes M_n(\bbC)^{\rm op}} \otimes (\partial_A \times \partial_B), 1_n \otimes 1_n^{\rm op} \otimes x]\| \\
        & = \|\gamma_{M_n(\bbC) \otimes M_n(\bbC)^{\rm op}} \otimes [\partial_A \times \partial_B, x]\| \\
        & = \|\gamma_{M_n(\bbC) \otimes M_n(\bbC)^{\rm op}}\| \|[\partial_A \times \partial_B, x]\| \\
        & = \| [\partial_A \times \partial_B, x]\| \\
        & = L_{\partial_A \times \partial_B}(x). \qedhere
    \end{align*}
\end{proof}

\end{document}
